\subsection{A Frobenius polynomial from ordinary neighbors}\label{subsec:frobenius-neighbors}

The Hecke action is now algebraic.  We next relate its reduction to
Frobenius by counting quotient isogenies at ordinary points of the
universal family.  This calculation depends on the reference
three-space and its Lie ranks; it is independent of the theta
multiplicity line.

For \(i=1,2,3\), let \(T_i\) be the full moduli Hecke operator with
\(i\) active exponents equal to one, all other selected exponents
zero, and multiplier exponent zero.  These are the active minuscule
cosets in the full similitude algebra.  They need not be positive in
\eqref{eq:positive-hecke-range}; Lemma~\ref{lem:hecke-albanese}
defines them by central clearing.

Let \(C_0\) be the central position with multiplier exponent one
and all active exponents zero.  On a nonactive selected block,
take all three exponents to be one if its Lie rank is three, and
zero if its Lie rank is zero.  Put
\begin{equation}\label{eq:integral-D-operators}
D_i=C_0^iT_i,\qquad 1\leq i\leq3.
\end{equation}
Here \(D_i\) has multiplier exponent \(i\), active exponents zero
or one, and nonactive selected exponents zero or \(i\).
It is a positive raw operator, hence an integral homomorphism of
\(A_S\).  The quotient branches defining \(D_i\) can therefore be
studied by specializing the isogenies constructed in
Subsection~\ref{subsec:algebraic-hecke}.

Let \(\widetilde A_S/\Fq\) be the good reduction at the fixed
completion of \(K\), and let \(\Phi\) be its \(q\)-power Frobenius
endomorphism.  The \(K\)-defined integral raw Hecke homomorphisms
extend over \(\Z_q\) by the Néron property; rational combinations
extend after tensoring the endomorphism ring with \(\Q\).
Their reductions are defined
over \(\Fq\), so they commute with \(\Phi\).  Use the same
letters \(D_i\) for these reductions.

\begin{proposition}\label{prop:albanese-congruence}
In the integral endomorphism ring of \(\widetilde A_S\),
\begin{equation}\label{eq:albanese-frobenius-polynomial}
\Phi^3-D_1\Phi^2+qD_2\Phi-q^3D_3=0.
\end{equation}
\end{proposition}

\begin{proof}
Let \(k=\overline{\Fq}\), and let \(x\) be an ordinary geometric
point of the special fiber of \(S\).  Its active \(q\)-divisible
group has the canonical splitting
\[
G_w=M_0\oplus E_0,
\]
where \(M_0\) is multiplicative of height one and \(E_0\) is
étale of height two.  Over the perfect field \(k\), the reduced
subgroup at each finite torsion level is the canonical étale
part; it is preserved by the endomorphism idempotents.  Each
nonactive selected block is wholly multiplicative or wholly
étale according to its Lie rank.

Three quotient constructions will account for the neighbors.
Write \(X_0\) for the quotient with active kernel \(M_0[q]\)
and multiplier \(q\), \(Y_0\) for the sum of quotients with
active kernel a line \(\ell\subset E_0[q]\) and multiplier
\(q\), and \(Z_0\) for the quotient with active kernel
\(E_0[q]\) and multiplier \(q^2\).  For multiplier \(q^h\),
prescribe the connected \(q^h\)-torsion on every nonactive
selected block and the forced annihilator on every opposite block.
We will prove that these quotients define full tuples of \(S\),
that \(X_0\) is Frobenius, and that the specialized neighbor sums are
\[
D_1:\ X_0+qY_0,\qquad
D_2:\ X_0Y_0+q^2Z_0,\qquad
D_3:\ X_0Z_0.
\]
The displayed formulas are the identities to establish; at this
point the quotient constructions have not yet been identified with
operations on \(S\).  Once they are, their cancellation will give
the asserted polynomial on Albanese differences.

\smallskip
\noindent\emph{Specialization multiplicities.}
Lift \(x\), using smoothness, to \(R=W(k)\), and write
\(\mathcal G_w\) for the lifted active \(q\)-divisible group.
The ring \(R\) is a complete henselian discrete valuation ring.
In the lifted active
\(q\)-torsion, the \(q^2\) open-and-closed special label
components lift uniquely by lifting idempotents of its finite
coordinate algebra.  Each lifted component is finite flat of
rank \(q\).  The identity-specializing component is a subgroup,
and multiplication respects the lifted labels by uniqueness.
Thus there is an exact sequence
\[
0\longrightarrow\mathcal H^0
\longrightarrow\mathcal G_w[q]
\longrightarrow(\Fq^2)_R\longrightarrow0
\]
with \(\mathcal H^0\) finite flat of rank \(q\).  No splitting
of this sequence over \(R\) is asserted.

Pass to a finite extension of \(R\) over which the generic torsion
is split.  Its active fiber is a three-dimensional \(\Fq\)-space
with distinguished line \(\mathcal H^0_{\mathrm{gen}}\).
For a fixed \(b\)-dimensional subspace of the label quotient,
the generic subspaces disjoint from that line and mapping onto
the label subspace are its graph lifts.  They form a torsor
under the \(b\)-dimensional space of maps to the line, so there
are \(q^b\) of them.  A generic subspace containing the line is
the unique full inverse image of its label subspace.

The schematic closure of a generic subgroup in the finite flat
torsion scheme is finite torsion-free over the new DVR, hence
flat.  It is still a subgroup, since the subgroup identities
extend from the generic fiber on flat schemes.  Intersection
with a label component is an idempotent direct summand of its
coordinate algebra, so its rank is its generic rank.  A graph
has rank one on each of its labels.  A finite flat rank-one
algebra over a local ring is that ring, and its special fiber
is the unique reduced point of that label component.  All
graph lifts over a label subspace therefore specialize to the
same étale subgroup of \(E_0[q]\).  A subspace containing the
distinguished line fills the corresponding rank-\(q\) components
and specializes to \(M_0[q]\) plus that étale subgroup.

The active kernels of \(D_i\) are \(i\)-dimensional subspaces of
the generic \(q\)-torsion.  The preceding count gives:
\[
\begin{array}{c|c|c}
\text{operator}&\text{special active kernel}&\text{multiplicity}\\ \hline
D_1&M_0[q]&1\\
   &\ell\subset E_0[q],\ \dim_{\Fq}\ell=1&q\text{ for each }\ell\\
D_2&M_0[q]\oplus\ell,\ \dim_{\Fq}\ell=1&1\text{ for each }\ell\\
   &E_0[q]&q^2\\
D_3&M_0[q]\oplus E_0[q]&1 .
\end{array}
\]
In the first two cases the total number is \(q^2+q+1\):
there are \(q+1\) étale lines.  The nonactive selected kernels
are uniquely their connected \(q^i\)-torsion for \(D_i\).

\smallskip
\noindent\emph{Full tuples and Frobenius.}
We now show that the counted kernels give targets in \(S\), and
identify the resulting quotient operations.  Properness of the
model of \(S\) extends each generic neighbor
over a finite valuation-ring extension.  The universal source
and target are abelian schemes over that new DVR.  Apply the
Néron mapping property there, anew, to extend the generic
isogeny \cite[\S1.2, Proposition~8]{BoschLutkebohmertRaynaud1990}.
The polarization identity extends by separatedness.  On each
fiber the kernel of the extension is contained in the finite
kernel of \(f^\vee\lambda'f=[q^i]\lambda\), so it has finite
kernel and is an isogeny between equal-dimensional abelian
varieties.  Since source and target polarizations have the
same degree, its degree is \(q^{3si}\) on every fiber.
Properness and fiberwise quasi-finiteness make the extension
finite.  Fiberwise flatness of isogenies between the smooth fibers
and the relative flatness criterion make it flat;
see \cite[\S7.3, Lemma~1 and Proposition~6; \S2.4,
Proposition~2]{BoschLutkebohmertRaynaud1990}.
Its finite flat kernel is therefore the
schematic closure of its generic kernel.

Finite-locally-free Cartier duality commutes with base change.
For a subgroup \(H\subset G\) in a finite group with a perfect
pairing, the annihilator is
\[
\operatorname{ann}(H)
 =\ker(G^D\longrightarrow H^D)=(G/H)^D.
\]
It is finite locally free and commutes with the above
specialization.  Thus the opposite generic kernels specialize
to the annihilators of the selected kernels just counted.
The corresponding nonactive kernels are already prescribed.
Equal total kernels give canonically isomorphic abelian
quotients.  Their descended \(\mathcal O_E\)-actions and
transported markings are unique, and their polarizations are
unique by cancelling \(f\) in \(f^\vee\lambda'f=q^i\lambda\).
The model retains these full tuples, so equal quotients are
equal target points of \(S\).  In particular every counted
kernel has such a target.  It is ordinary, since ordinariness
is invariant under isogeny.

The constructions \(X_0,Y_0,Z_0\) are therefore operations on
the ordinary points of the same full moduli space.  We use them
as multiset operations on points; no global special-fiber
algebraic correspondence is required.

The operation \(X_0\) is the \(q\)-power Frobenius on full tuples.
Indeed its total kernel is the connected \(q\)-torsion on every
block.  On the active pair, duality pairs the multiplicative
part with the opposite étale part and the étale part with the
opposite multiplicative part, so the annihilator of \(M_0[q]\)
is exactly the opposite connected part; the pure nonactive
pairs have the same property.  The total kernel is therefore
the kernel of relative Frobenius
\(F_B:B\to B^{(q)}\).  Naturality of the polarization and
Frobenius--Verschiebung duality give
\[
F_B^\vee\lambda^{(q)}F_B
 =V_{B^\vee}F_{B^\vee}\lambda=[q]\lambda.
\]
Here \(V_{B^\vee}\) denotes Verschiebung.
The action on the target is \(\iota_E^{(q)}\).  For the constant
label group \(T=(\Z/N\Z)^{6s}\), naturality gives
\[
F_B\eta_N=\eta_N^{(q)}F_T.
\]
Relative Frobenius on \(T\) fixes its labels, and \(F_B\) is an
isomorphism on \(N\)-torsion.  Thus the transported marking is
exactly the twisted marking, with no multiplication-by-\(q\)
marking change.  Since the model and universal data are defined
over \(\Fq\), this is its coordinate-\(q\)-power operation on
geometric points, not its inverse.

Frobenius commutes with \(Y_0\) and with \(Z_0\) as multiset
operations.  For any of their quotient maps \(f:B\to B'\),
\[
F_{B'}f=f^{(q)}F_B.
\]
Frobenius is an isomorphism on the étale part and permutes its
line choices bijectively; this square also respects the
polarizations, actions, and markings.  On the active block the
composite \(X_0Y_0\) has kernel \(M_0[q]\oplus\ell\) with
multiplier \(q^2\), and \(X_0Z_0\) has kernel
\(M_0[q]\oplus E_0[q]\) with multiplier \(q^3\).

For clarity, the opposite kernels of these composites are
indeed forced by their selected kernels and multipliers.  On
the \(q\)-divisible groups, write \(G,G'\) for source and target
and define the polarization adjoint
\(f^\dagger=\lambda^{-1}f^D\lambda'\); the polarizations are
isomorphisms there because their degrees are prime to \(q\).
For a composite with multiplier \(q^h\),
\[
f^\dagger f=[q^h],\qquad ff^\dagger=[q^h].
\]
The second identity follows by cancellation of the fppf
epimorphism \(f\).  As fppf subgroup sheaves,
\(\ker f=f^\dagger(G'[q^h])\).  For \(u,v\in G'[q^h]\),
adjunction gives
\[
e_{q^h}\bigl(f^\dagger u,f^\dagger v\bigr)
 =e'_{q^h}\bigl(u,ff^\dagger v\bigr)=1,
\]
so the image is isotropic in \(G[q^h]\).  The endomorphism
idempotents split its kernel into conjugate-block parts.
If \(K_j,K_{\bar j}\) are the parts in one conjugate pair of
height-three blocks, then
\[
K_{\bar j}\subseteq K_j^\perp,\qquad
\operatorname{rank}(K_j)\operatorname{rank}(K_{\bar j})
 \leq q^{3h}.
\]
Cartier duality gives equal kernel ranks for \(f\) and \(f^\dagger\).
Their composite has kernel rank \(q^{6sh}\), so
\(\operatorname{rank}\ker f=q^{3sh}\).
Multiplying the pairwise inequalities gives equality, and hence
equality in each pair.  Thus \(K_{\bar j}=K_j^\perp\).
For \(X_0Z_0\), for example, the annihilator of the active
\(q\)-torsion in \(q^3\)-torsion is the opposite block's
\(q^2\)-torsion.  The pure nonactive connected kernels likewise
compose to connected \(q^h\)-torsion.  Uniqueness of the
polarizations and markings now identifies the composite targets
with the counted targets.  We have proved the multiset identities
\begin{equation}\label{eq:ordinary-neighbor-counts}
\begin{aligned}
D_1&:\quad X_0+qY_0,\\
D_2&:\quad X_0Y_0+q^2Z_0,\\
D_3&:\quad X_0Z_0.
\end{aligned}
\end{equation}

\smallskip
\noindent\emph{The integral Albanese identity.}
The neighbor counts are now identities of multisets of full tuples.
We pass from them to endomorphisms by applying the Albanese difference
formula.  Let \(\mathcal O\) be the set of ordinary \(k\)-points of the
union of components and extend \(X_0,Y_0,Z_0\) linearly to the
free abelian group \(\Z[\mathcal O]\).  Define
\[
\alpha:\Z[\mathcal O]\longrightarrow\widetilde A_S(k),
\qquad \alpha([x])=a_{j(x)}(x).
\]
The component Albanese maps reduce over \(\Fq\), so
\(\Phi\alpha=\alpha X_0\).  For ordinary \(x,y\) in one source
component, lift them and specialize the integral identity
\eqref{eq:hecke-albanese-differences}.  With
\[
N_1=X_0+qY_0,\qquad
N_2=X_0Y_0+q^2Z_0,\qquad
N_3=X_0Z_0,
\]
it gives
\[
D_i\bigl(a_j(x)-a_j(y)\bigr)
 =\alpha\bigl(N_i([x]-[y])\bigr).
\]
This is an equality in the actual group of geometric points,
with all branch multiplicities retained.  No rationalization
of a finite-field point group is used.

Frobenius preserves each component and its ordinary locus.
Applying the same formula to \(X_0^m x,X_0^m y\) shows that
the endomorphism in \eqref{eq:albanese-frobenius-polynomial}
on this difference is
\[
\alpha\!\left(
(X_0^3-N_1X_0^2+qN_2X_0-q^3N_3)([x]-[y])
\right).
\]
The operator inside this expression vanishes:
\[
\begin{aligned}
&X_0^3-N_1X_0^2+qN_2X_0-q^3N_3\\
&\qquad=-qY_0X_0^2+qX_0Y_0X_0
        +q^3Z_0X_0-q^3X_0Z_0=0.
\end{aligned}
\]
Only the two proved commutations with \(X_0\) are used.
Ordinary points are dense in every smooth component by
Lemma~\ref{lem:good-ordinary-primes}.  Therefore the morphism
obtained by applying the endomorphism to
\(a_j(x)-a_j(y)\) vanishes on all of the product of that
component with itself.  The finite-sum generation in the same
lemma makes the endomorphism zero on each Albanese factor.
It is zero on \(\widetilde A_S\) after base change to \(k\);
faithful base change gives the integral identity over \(\Fq\).
\end{proof}

For the cohomological application, we record the Frobenius convention.  On
\(H^1_{\mathrm{cris}}(\widetilde A_S/\Z_q)[1/q]\), crystalline
Frobenius is the pullback \(\varphi=\Phi^*\) of this \(q\)-power
endomorphism: it is defined by absolute-Frobenius pullback, and
the Witt Frobenius on \(W(\Fq)=\Z_q\) is the identity because
\(q\) is prime \cite[Tag~07N0]{StacksProject}.  In degree-one
cohomology pullback is additive in abelian homomorphisms.
It reverses composition, but the displayed endomorphisms
commute.  Thus \eqref{eq:albanese-frobenius-polynomial} gives
the same operator polynomial with \(\varphi\) and the Hecke
pullbacks, not a reciprocal polynomial.  This is the identity
used in Proposition~\ref{prop:active-hodge-types}.
